\section{无线实验平台的选择与比较}
\label{sec:2}

系统比较不同无线配置对端到端训练流程的影响，对实验平台的收发能力提出了明确的技术诉求。一方面，参数与发送行为必须可控，从而支持研究人员主动构造差异化的实验条件，在保持模型结构、训练数据和本地计算负载完全一致的前提下，观察通信耗时与轮次推进的变化。另一方面，平台应当支持高效的一对多模型分发。联邦学习每轮下行都需要将完全相同的全局模型推送到全体参与客户端，广播传输可以有效避免单播方式带来的重复发送开销。尽管单播同样能够完成真实的无线联邦学习实验，但随着客户端数量增加，广播在减少网络占用方面的优势会更加明显。

围绕上述实验目标，研究者通常有几种不同的实验平台搭建路径，各自存在不同的工程特征与使用条件（见 Table~\ref{tab:1}）。常见的做法是采用消费级 Wi-Fi 的常规联网方式。商用无线网卡与接入点容易获得且价格低廉，网络配置相当成熟。但常规网络协议栈把物理层收发细节封装在固件与驱动之内，固件内置的速率自适应机制会在后台频繁升降调制阶数，导致实验人员很难锁定一组预设的无线参数进行受控对比。同时，标准接入点往往将广播帧限制在较低的基本速率集，分发数十兆字节的模型文件需要消耗较多时间。第二种途径是基于低功耗蓝牙构建测试环境。低功耗蓝牙硬件小巧、能耗极低且单价只有数十元，非常适合大规模低功耗节点的互联。但在协议标准规范下，其物理层数据速率通常在 1 至 2 Mbps 之间，单包有效载荷也较小，面对数十兆字节的深度学习模型文件时分发耗时漫长，更适合承载参数量较小的轻量级模型或局部更新。若追求更高的物理层灵活性，研究者往往转向软件定义无线电。通过软件定义无线电平台，研究人员可以直接在基带层面灵活定制调制格式、纠错编码与射频参数，具备极高的可编程性。然而，这类硬件的采购成本显著偏高，单节点设备从数千元到数万元不等。更为关键的是，研究团队需要自行实现或深度移植完整的介质访问控制与高层传输协议栈，系统维护工作繁重，且在部分频段发射无线信号时还需要取得专门的无线电管理频谱许可。

% 原稿标记保留：**[Table 1 占位 无线实验平台搭建方案对比矩阵]**

% 原稿题注保留：> Table 1 按关键无线参数可控性、广播分发能力、硬件成本与部署维护条件四个维度，对比四种无线实验搭建方案的实现特征与使用条件。

\begin{table*}[tb]
\caption{无线实验平台搭建方案对比。}
\label{tab:1}
\centering
\small
\renewcommand{\arraystretch}{1.25}
\setlength{\tabcolsep}{5pt}
\begin{tabularx}{\textwidth}{@{}>{\raggedright\arraybackslash}p{0.12\textwidth}YYYY@{}}
\toprule
\textbf{评估维度} & \textbf{采用消费级 Wi-Fi 常规联网方式的方案} & \textbf{基于低功耗蓝牙的方案} & \textbf{基于软件定义无线电的定制方案} & \textbf{本文平台（WFB-FL）} \\
\midrule
\textbf{关键无线参数可控性} & 受固件自适应调节与驱动封装限制，通常难以在运行中锁定发送参数 & 参数主要由协议栈固化配置管理，物理层可调范围较窄 & 支持基带与射频参数的灵活软件定制 & 支持显式指定信道、带宽、MCS、保护间隔、FEC 与发射功率 \\
\textbf{广播分发能力} & 通常受限于接入点协商的基本速率集，亦可采用单播替代广播 & 受协议标准速率（1 至 2 Mbps）约束，分发大体积模型文件耗时较长 & 具备物理广播能力，但需自行实现可靠分发与重传协议栈 & 支持以高阶速率发送广播帧，由文件分发协议处理可靠交付 \\
\textbf{硬件成本} & 较低（单节点百元级商用网卡或路由设备） & 低（单节点数十元级射频模块） & 较高（单节点数千元至数万元不等） & 较低（单节点百元级消费级 USB 网卡） \\
\textbf{部署与维护条件} & 常规网络配置成熟，接入点部署方便 & 拓扑配置简单，管理开销小 & 需基带软硬件联合调试，在受保护频段需无线电管理许可 & 工作于免许可频段，需编译补丁驱动并调试收发协同逻辑 \\
\bottomrule
\end{tabularx}
\end{table*}

面向本文受控对比的实验目标，WFB-FL 将消费级 802.11 硬件、可设置的收发参数与面向模型分发的广播传输组合在同一方案中。平台依托双频 USB 无线网卡，利用开源驱动补丁开放监听模式（monitor 模式）与原始帧注入能力。在免许可频段上，发送端绕过标准连接协商，直接构造包含预设信道、带宽、MCS、保护间隔、FEC 冗余度与发射功率的物理帧，并以高阶速率向多客户端广播模型文件；通信对端则捕获空中实际到达的无线帧并记录元数据。这种机制为受控实验提供了主动调整物理参数与观察实际通信表现的途径，但也对应着具体的部署与维护代价。这套系统并非开箱即用。研究人员需要选用支持注入模式的特定芯片网卡，在各边缘节点上独立编译专用驱动补丁，并手动配置免许可频段的工作模式与节点间的调度时序。百元级的消费级硬件提供了切实的设备可获得性，而收发控制与多节点防碰撞则需要完全交由软件逻辑独立承担。
